% Français 9115, batch 35 — Gallica views 681–700.
% Pass: Opus 5.5 (claude-opus-5-5), 2026-09-22 — first pass, unchecked against the views by a human.
%
% What the twenty views hold. Heat, as in batch 34: Fourier's problem of
% the infinite solid cylinder whose sections are heated alike, then the
% opening of a piece on the sphere, the cone and the cylinder. Nine
% views carry text, eleven do not. No number theory, no microfilm target,
% no printed matter, no repeated exposure. Single leaves hinged on grey
% mounts, alternately high and low; the ink series stands on the recto.
%
%   v. 682–691  the end of the cylinder piece begun at v. 672 (batch 34),
%               her pages 6–10 (page 5 is v. 680). v. 682 (ink 335) picks
%               up the surface condition hv + dv/dr = 0 at r = R: the
%               determined equation in g, reality of its roots « N.o 306
%               et suivans »; the general integral as a sum over β of
%               (A^β sin βφ + B^β cos βφ) Σ a e^{−gkt} s (her art. 4).
%               v. 684 (ink 336): the roots, values and coefficients
%               named (the lines struck on the back of v. 678); N.o 5, the
%               series for s rewritten in θ, the differential equations
%               of successive orders, art. 308. v. 686 (ink 337): the roots
%               are real; N.o 6, the factor A^β sin βφ + B^β cos βφ, β
%               infinite and β null two ways of presenting one state; the
%               g^β grow with β. v. 688 (ink 338): for large t only the
%               first two orders remain; two points at the ends of a
%               diameter. v. 690 (ink 339): their half-sum is the mean
%               temperature of the layer; the points at π/2 from the
%               extremes cool as in a cylinder heated uniformly; the
%               analogy with the ring, « art 245 p 277 dela Théorie dela
%               chaleur ». v. 691, its back: five lines closing the piece
%               (each section behaves as a round plate shielded from the
%               air; the ring is the case of constant radius). A draft:
%               much corrected between the lines on v. 686–688.
%   v. 692–693  the title « Equation du mouvement varié dela chaleur dans
%               un cylindre solide » and articles 1–3 written a second
%               time, almost word for word as at v. 672–676, in a steady
%               hand without corrections; v. 692 ink 340, no page number;
%               v. 693, its back, page 2; stops on « Pour y parvenir il ».
%   v. 694–695  pages 7 and 8 of that second copy (ink 341 on 694, 695 its
%               back): the text of the draft's v. 686–691 from « manières
%               differentes de présenter » to the end, recopied almost
%               verbatim. Its pages 3–6 are not in this batch; nothing on
%               the leaves says which of the two copies was written first.
%   v. 696      (ink 342, her page 1) a new piece, heavily corrected
%               draft: « Equations du mouvement dela chaleur dans la
%               sphère, le cône et le cylindre solides échauffés d'une
%               manière entièrement arbitraire »: Laplace and Poisson
%               transform the general equation by a change of coordinates,
%               which hides the connection between the three equations;
%               with « la méthode employée par M.r fourier » each case is
%               examined specially, and the molecule of the sphere is seen
%               to belong to a cone with vertex at the centre. Stops on
%               « le même coté »; back blank (697).
%   v. 698      (ink 343) the sphere heated uniformly from centre to
%               circumference: the angular terms vanish, no heat passes
%               between the quadrilaterals cut by meridians and parallels;
%               the equation of their vanishing read as a balance of heat
%               through the meridian and parallel arcs. Back blank (699).
%   v. 700      a leaf numbered 344 on its blank side; its text is on the
%               back, v. 701 (outside the batch).
%
% Views not transcribed. v. 681, 683, 685, 687, 689: blank backs of the
% leaves 334–338, the recto showing through reversed (checked at band
% scale). v. 697, 699: blank backs of 342, 343. v. 700: blank but for the
% ink number 344, recorded in a note at v. 698.
%
% Seams. At 680/681: v. 680 (batch 34) ends « qui doit être satisfaite
% lorsque » and v. 682 opens « $r$ a la valeur totale $R$; on en conclura
% l'équation déterminée »: the sentence runs over the blank back v. 681.
% Her page 5 at v. 680 (read doubtfully in batch 34) is confirmed by the
% run 6–10 at v. 682–690. The lines struck on v. 679 reappear, unstruck, at
% the head of v. 684. At 700/701: v. 700 is the numbered blank side of leaf
% 344 and v. 701 (glanced at) its written back, on the sphere (« … par le
% pôle dela sphère on mène un nombre infini de méridiens … »); v. 698 ends
% its paragraph on « θ − dθ: », so whether 701 follows it directly is for
% batch 36 to say.
%
% The numbers on the leaves, and why no \folio{} is written. (1) The ink
% series at the top right, in the larger rounder hand, one per leaf on
% the side that faces the camera first: 335 (v. 682), 336 (684), 337
% (686), 338 (688), 339 (690), 340 (692), 341 (694), 342 (696), 343 (698),
% 344 (700, on the blank side of its leaf); 334 on v. 680 just before.
% None on the backs 691, 693, 695. Penned, as throughout the volume: no
% \folio{}. (2) Her own page numbers, small, beside or above the ink
% number and struck through with an oblique stroke: 6 (682), 7 (684), 8
% (686), 9 (688), 10 (690) for the draft; none on 692, « 2 » underlined
% at the top right of its back (693, not struck), « 7 » struck (694),
% « 8 » not struck (695) for the second copy; « 1 » struck (696) for the
% sphere piece. (3) Her article numbers N.o 5 (684), N.o 6 (686), 1, 2
% (692), 3 (693); references to her art 4 (688, 694) and art 5 (686,
% 694); to « N.o 306 et suivans », « l'art 308 », « art 245 p 277 dela
% Théorie dela chaleur » (682, 684, 690, 695). Every number was read on
% its view; none is inferred.
%
% Manuscript C. Not here. Nothing in these views is number theory, and
% the ink series runs 335–344. If Laubenbacher and Pengelley's ff.
% 348r–349r are this series, they would fall about four leaves on, near
% views 708–710 at one numbered leaf per two views: an extrapolation, not
% a reading.
%
% Notation. Hers, kept. β is her looped letter with a stem; the
% superscripts ′, 2, β on g, A, B, a, s, u are order indices, not powers,
% and are set as she writes them; subscripts 1, 2, 3 likewise. K, C.D,
% h, k, g as in batch 34. « sin », « cos », « tang » with or without a
% stop, set \text{sin} etc. Her « &c » is \&c. The first series of v. 684
% is split over two rows, which the leaf does not do. Letters she underlines in
% the prose are \underline inside maths; her underlined words would be
% \emph{}. Plurals cancelled by a small cross are given in the singular,
% with a note (v. 696). The long s is transcribed s. Spelling hers, not
% flagged: « tems », « mêmetems », « suivans », « differens »,
% « coëfficiens », « exposans », « appartenans », « parconséquent »,
% « cequi », « cequ'on », « dela », « delà », « aulieu », « alafois »,
% « lamême », « laquantité », « laquelle », « invelappe », « apperçoit »,
% « appercevoir », « paralelle » beside « parallèle », « les point »,
% « reçus », « est correspond » (v. 691). Where the leaf and the
% calculation disagree the leaf stands, with a note (v. 682, 684, 690,
% 695, 698).
%
% What could not be read, and what is offered with doubt. \ill{} stands
% four times: a small added word at the right edge of v. 688; a word
% begun after « les trois équations. » at v. 696; one word of a fine
% interlinear addition and one struck word under a blot at v. 696.
% \uncertain{} stands for: « T. 2 » and the exponent 0 of B (v. 682);
% f before θ, the exponent 3, and the exponent 0 of B (v. 684);
% « indiqueroit » (v. 686); « sensiblement » (v. 688); v in the first
% equation (v. 692); the ending of « donnent », the interlinear « qu'on
% cherche la … qu'elle et », and the ending of « détermine » (v. 696);
% v in dv/dt (v. 698).
%
% Nothing here was checked against any published edition of Sophie
% Germain's papers, or against Fourier's Théorie analytique de la
% chaleur, Laplace or Poisson.
\documentclass[11pt,a4paper]{article}
% The preamble shared by every transcription of the mathematicians' manuscripts in Gallica.
%
% It rests on one idea: **uncertainty compounds, it does not resolve.** A
% transcription that smooths over a doubtful word has destroyed the only thing
% it was made for — a reader must be able to tell, at every point, what was on
% the page and what was supplied. The apparatus macros below are therefore the
% critical apparatus, not decoration.
%
% The same commands are recognised by `scripts/render.mjs`, which produces the
% site's reading view, and by `scripts/tei.mjs`. A macro added here must be
% added there too, or rendering fails loudly — which is the wanted behaviour.
%
% Comments here are English, like the rest of the repository. The documents
% this preamble sets are French, and that is deliberate: see the skills.

\ProvidesPackage{germain}[2026/09/15 Transcriptions of mathematicians' manuscripts in Gallica]

\usepackage{amsmath,amssymb,amsthm}
\usepackage{mathrsfs}
% Kept from the parent preamble so the renderer's diagram support stays
% available; these manuscripts draw no commutative diagrams, and nothing here uses it.
\usepackage{tikz-cd}
\usepackage[english,french]{babel}
\usepackage{fontspec}

% --- Greek ------------------------------------------------------------------
% The text face stays fontspec's default, Latin Modern, which has no Greek:
% XeTeX drops every glyph it lacks with a line in the log and nothing on the
% page, and the Greek words the manuscripts quote vanished from the PDFs. So
% the Greek and Greek Extended blocks get a XeTeX character class of their own,
% and entering or leaving a run of it switches the family to CMU Serif — the
% same Computer Modern design, with polytonic Greek — and back. Only the
% family changes: a Greek word in \textit or \textbf keeps its shape and series.
% The transitions to and from every other class carry the tokens that class
% has towards an ordinary letter, so French spacing before « ; » or inside
% « » still holds after a Greek word. No grouping, so no brace or \endgroup
% can come out unbalanced; maths is left alone, since XeTeX inserts nothing
% there. Kept identical in `exercices.sty`.
\makeatletter
\newfontfamily\germainGreekfont{cmun}[
  Extension=.otf, NFSSFamily=germaingreek,
  UprightFont=*rm, ItalicFont=*ti, SlantedFont=*sl,
  BoldFont=*bx, BoldItalicFont=*bi, BoldSlantedFont=*bl]
\newXeTeXintercharclass\germain@greek
\def\germain@greekrange#1#2{%
  \@tempcnta=#1\relax
  \loop
    \XeTeXcharclass\@tempcnta=\germain@greek
    \ifnum\@tempcnta<#2\advance\@tempcnta\@ne\repeat}
\germain@greekrange{"0370}{"03FF}
\germain@greekrange{"1F00}{"1FFF}
\def\germain@greekprev{\rmdefault}
\def\germain@greekin{%
  \edef\germain@greekprev{\f@family}\fontfamily{germaingreek}\selectfont}
\def\germain@greekout{\fontfamily{\germain@greekprev}\selectfont}
\def\germain@greekjoin{%
  \XeTeXinterchartokenstate=\@ne
  \@tempcnta=\z@
  \loop
    \ifnum\@tempcnta=\germain@greek\else
      \XeTeXinterchartoks\germain@greek\@tempcnta=\expandafter{%
        \expandafter\germain@greekout\the\XeTeXinterchartoks\z@\@tempcnta}%
      \XeTeXinterchartoks\@tempcnta\germain@greek=\expandafter{%
        \the\XeTeXinterchartoks\@tempcnta\z@\germain@greekin}%
    \fi
    \ifnum\@tempcnta<4095\advance\@tempcnta\@ne\repeat}
% babel-french sets its own classes, and saves and restores the interchar
% state, each time a language is selected: join again after it does.
\addto\extrasfrench{\germain@greekjoin}
\addto\extrasenglish{\germain@greekjoin}
\AtBeginDocument{\germain@greekjoin}
\makeatother

\usepackage{geometry}
\geometry{a4paper,margin=25mm,marginparwidth=28mm}
\usepackage{marginnote}
\usepackage{xcolor}
\usepackage[normalem]{ulem}
\setlength{\emergencystretch}{3em}
\usepackage{hyperref}
\hypersetup{colorlinks=true,linkcolor=black,urlcolor=black,pdfborder={0 0 0}}

% --- Batch metadata ---------------------------------------------------------
% Taken as they stand by the rendering script, which shows them at the head of
% the reading view. They fix what the file claims to cover. The macro names are
% the parent project's, so that the scripts stay shared; \folder here means the
% volume's slug — `fr-9115` — and \pages counts Gallica views.

\newcommand{\folder}[1]{\gdef\grothFolder{#1}}
\newcommand{\batch}[1]{\gdef\grothBatch{#1}}
\newcommand{\pages}[2]{\gdef\grothFirst{#1}\gdef\grothLast{#2}}
\newcommand{\foldertitle}[1]{\gdef\grothFoldertitle{#1}}

% The holder's own shelfmark, printed as they write it: « Français 9115 ».
\newcommand{\shelfmark}[1]{\gdef\grothShelfmark{#1}}

% The dating, copied verbatim from the holder's catalogue — never inferred
% here. « XIXe siècle » is the BnF's; a pass that dates a leaf from its content
% says so in a \note{}, as its own inference.
\newcommand{\dating}[1]{\gdef\grothDating{#1}}

% The status notice, printed in the title block and repeated as a diagonal
% watermark on every page of the PDF. The manuscripts are in the public
% domain; what this line declares is not a right but a quality — an
% unverified first pass by a machine. The skills write the line; a file
% without it gets no watermark, so leaving it out is visible in the output.
\newcommand{\watermark}[1]{\gdef\grothWatermark{#1}}

\gdef\grothFolder{?}\gdef\grothBatch{}\gdef\grothFirst{?}\gdef\grothLast{?}%
\gdef\grothFoldertitle{}\gdef\grothShelfmark{}\gdef\grothDating{}\gdef\grothWatermark{}

% --- The résumé -------------------------------------------------------------
\newenvironment{resume}
  {\small\setlength{\parskip}{0.45em}}
  {\par\medskip\hrule\medskip}

% --- Keywords ---------------------------------------------------------------
% In English, closing the modernised reading's résumé: the modern vocabulary
% under which to search for what the leaves construct. The manifest extracts
% them to tag the volume — this line is the single source of the tags.
\newcommand{\keywords}[1]{\par\smallskip\noindent
  {\footnotesize\textsc{Keywords} --- \textit{#1}}\par}

% --- The critical apparatus -------------------------------------------------

% The Gallica view number, in the margin. This is the anchor that lets a
% disagreement over a formula be settled by looking at *one* image rather than
% a volume of 750. The site uses it to turn the facsimile as the transcription
% is scrolled. A view is one scanned image: usually one side of a leaf.
\newcommand{\page}[1]{%
  \marginnote{\footnotesize\color{gray!70}\textsf{v.\,#1}}%
  \label{page:#1}\ignorespaces}

% The BnF's pencilled foliation, where the leaf carries it and a pass can read
% it: « 198r ». This is what the literature cites — « ff. 198r–208v » — and
% the one line that turns a citation into a view. Recorded, never inferred:
% a folio a pass cannot read is not written.
\newcommand{\folio}[1]{%
  \marginnote{\footnotesize\color{gray!50}\textsf{f.\,#1}}\ignorespaces}

% The modernised reading's coarser counterpart to \page{}: which views a
% section is drawn from.
\newcommand{\pagerange}[2]{%
  \marginnote{\footnotesize\color{gray!70}\textsf{v.\,#1--#2}}%
  \ignorespaces}

% Illegible: never guessed.
\newcommand{\ill}{\textcolor{red!60!black}{[\,\ldots\,]}}

% A doubtful reading: the word is given, and flagged as uncertain.
\newcommand{\uncertain}[1]{\uline{#1}}

% An editorial addition — a missing letter, an implied word.
\newcommand{\add}[1]{\textcolor{blue!50!black}{[#1]}}

% Struck out by the writer. What was crossed out is part of the document.
\newcommand{\struck}[1]{\sout{#1}}

% The transcriber's note, on the state of the page or the mathematics.
\newcommand{\note}[1]{\par\smallskip{\small\color{gray!60!black}\textit{[#1]}}\par\smallskip}

% A marginal note of hers, distinct from the body of the text.
\newcommand{\marginal}[1]{\par\smallskip\noindent\hspace{1em}%
  {\small\color{gray!60!black}#1}\par\smallskip}

% --- Title ------------------------------------------------------------------
% The title block is French, like the documents themselves.

\AddToHook{shipout/background}{%
  \ifx\grothWatermark\empty\else
    \begin{tikzpicture}[remember picture, overlay]
      \node[rotate=52, scale=3.4, align=center, text=gray!18,
            font=\sffamily\bfseries]
        at (current page.center) {\grothWatermark};
    \end{tikzpicture}%
  \fi}

\AtBeginDocument{%
  \begin{center}
    {\large\bfseries Manuscrits de mathématiciens (Gallica) ---
      \ifx\grothShelfmark\empty\grothFolder\else\grothShelfmark\fi}\\[2pt]
    {\itshape\grothFoldertitle}\\[4pt]
    {\small vues \grothFirst--\grothLast
      \ifx\grothBatch\empty\else\ (lot \grothBatch)\fi}\\[2pt]
    \ifx\grothDating\empty\else
      {\small\color{gray!60!black}Datation du catalogue : \grothDating}\\[2pt]
    \fi
    \ifx\grothWatermark\empty\else
      {\small\bfseries\color{red!45!gray}\grothWatermark}\\[2pt]
    \fi
    {\footnotesize\color{gray!60!black}Transcription établie d'après les images IIIF de Gallica.
     Source gallica.bnf.fr / Bibliothèque nationale de France.\\
     Les lectures incertaines sont soulignées, les illisibles notées [\,\ldots\,],
     les ajouts entre crochets ; \textsf{v.} numérote les vues de Gallica,
     \textsf{f.} la foliotation au crayon de la BnF.}
  \end{center}
  \vspace{1em}}

\endinput


\folder{fr-9115}
\shelfmark{Français 9115}
\batch{35}
\pages{681}{700}
\dating{1801-1900}
\watermark{Lecture automatique\\non vérifiée}
\foldertitle{Papiers de Sophie GERMAIN. — Recueil de dissertations et problèmes mathématiques et physiques.}

\begin{document}

\note{Numérotation des feuillets. En haut à droite des feuillets, à
l'encre, d'une écriture plus ronde, un nombre de la série qui court dans
tout le volume : 335 (vue 682), 336 (684), 337 (686), 338 (688), 339
(690), 340 (692), 341 (694), 342 (696), 343 (698), 344 (700). Cette série
occupe la place de la foliotation de la Bibliothèque, mais elle est
tracée à la plume et non au crayon, et aucun « f. » n'est donc porté.
Près de ce nombre, sa propre pagination, petite, le plus souvent barrée
(voir les notes de chaque vue). Tous ces nombres ont été lus sur la vue ;
aucun n'est déduit. Les vues 681, 683, 685, 687, 689, 697 et 699 sont
des dos blancs où l'écriture du recto se voit à l'envers ; la vue 700 est
le côté blanc du feuillet 344, dont le texte est au dos (vue 701, hors
de ce lot). Elles ne sont pas transcrites.}

\page{682}

$\underline{r}$ a la valeur totale $R$; on en conclura l'équation
déterminée:
\[ b\left(1 - \frac{gR^2}{2(n+1)} + \frac{g^2R^4}{2.4(n+1)(n+3)} - \frac{g^3R^6}{2.4.6(n+1)(n+3)(n+5)} + \&c\right) \]
\[ = \frac{2gR}{2(n+1)} - \frac{4g^2R^3}{2.4(n+1)(n+3)} + \frac{6g^3R^6}{2.4.6(n+1)(n+3)(n+5)} - \&c \]
$g$, $g'$, $g^2 \ldots g^\beta$. étant cequi devient laquantité $g$ lorsque
\struck{que} le coëfficient de $\varphi$ est $0, 1, 2 \ldots \beta$. on prouvera,
en se conformant a cequi est demontré N.\textsuperscript{o} 306 et suivans \emph{\uncertain{T. 2}}
que quelque soit le nombre $\underline{n}$ cette equation en
$g$, $g'$, $g^2 \ldots g^\beta$, dans laquelle $\underline{b}$ et $R$ sont des quantités
déterminées a une infinité de racines réelles.

Il suit delà qu'on peut donner a la variable $\underline{v}$ une infinité
de valeurs des formes. $e^{-gkt}s$, $e^{-g'kt}(A'\,\text{sin}\,\varphi + B'\,\text{cos}\,\varphi)\,r\,s$,
$e^{-g^2kt}(A^2\,\text{sin}.2\varphi + B^2\,\text{cos}\,2\varphi)\,r^2s \ldots\ldots e^{-g^\beta kt}(A^\beta\,\text{sin}\,\beta\varphi + B^\beta\,\text{cos}\,\beta\varphi)\,r^\beta s$.
qui differeront seulement par les exposans $g$, $g'$, $g^2 \ldots g^\beta$.
On pourra donc composer une valeur plus générale en ajoutant
toutes ces valeurs \struck{multipliées} particulières multipliées par des coëfficiens
arbitraires.

L'intégrale qui servira a résoudre dans toute son étendue
la question proposée est donnée par l'équation:
\[ v = B^{\uncertain{0}}\left(a_1e^{-g_1kt}s_1 + a_2e^{-g_2kt}s_2 + a_3e^{-g_3kt}s_3 + \&c\right) \]
\[ + (A'\,\text{sin}\,\varphi + B'\,\text{cos}\,\varphi)\left\{a'_1e^{-g'_1kt}s'_1 + a'_2e^{-g'_2kt}s'_2 + a'_3e^{-g'_3kt}s'_3 + \&c\right\} \]
\[ + (A^2\,\text{sin}\,2\varphi + B^2\,\text{cos}\,2\varphi)\left\{a^2_1e^{-g^2_1kt}s^2_1 + a^2_2e^{-g^2_2kt}s^2_2 + a^2_3e^{-g^2_3kt}s^2_3 + \&c\right\} \]
\&c.
\[ + (A^\beta\,\text{sin}\,\beta\varphi + B^\beta\,\text{cos}\,\beta\varphi)\left\{a^\beta_1e^{-g^\beta_1kt}s^\beta_1 + a^\beta_2e^{-g^\beta_2kt}s^\beta_2 + a^\beta_3e^{-g^\beta_3kt}s^\beta_3 + \&c\right\} \]
\note{En haut à droite, à l'encre, « 335 », le 5 à longue queue ; au-dessus,
un « 6 » barré d'un trait oblique, de sa main : sa pagination du morceau
(voir l'en-tête). Au début de la première ligne, « a » est repassé
d'une encre plus lourde sur une lettre qu'il couvre, peut-être une
majuscule. Le troisième terme du second membre porte « $R^6$ » sur la
feuille, là où le calcul demande $R^5$ ; on garde la feuille. Dans les
exposants et les indices, la lettre bouclée à hampe est lue $\beta$, comme
dans les pages précédentes du volume ; les exposants $'$, $2$, $\beta$
de $g$, $A$, $B$, $a$, $s$ sont des indices d'ordre, non des puissances.
L'exposant de $B$ dans la première ligne de l'intégrale est un pâté ;
« 0 » est offert d'après la vue 684, qui écrit $A^0$, $B^0$. Les
deux premiers $s$ de cette ligne sont chargés d'encre. Le « que » en tête de la ligne qui suit « lorsque » paraît
biffé d'un trait bouclé. Après « suivans », un petit groupe
souligné dont la fin est chargée d'encre, lu avec doute « T. 2 ». Le bas du feuillet est blanc.}

\page{684}

$g_1, g_2, g_3\,\&c.\ g'_1, g'_2, g'_3\,\&c.\ g^2_1, g^2_2, g^2_3\,\&c. \ldots\ g^\beta_1, g^\beta_2, g^\beta_3\,\&c.$
\struck{qui} désignent toutes les valeurs de $g$, $g'$, $g^2 \ldots g^\beta$, qui satisfont
aux équations déterminées correspondantes aux valeurs $1, 3, 5 \ldots 2\beta+1$
du nombre $n$; $u_1, u_2, u_3\,\&c.\ u'_1, u'_2, u'_3\,\&c.\ u^2_1, u^2_2, u^2_3\,\&c. \ldots$
$u^\beta_1, u^\beta_2, u^\beta_3\,\&c.$ désignent les valeurs de $\underline{u}$ qui correspondent à ces
differentes racines; $B^{\uncertain{0}}$; $A'$, $B'$; $A^2$, $B^2$; $\ldots$ $A^\beta$, $B^\beta$; et $a_1, a_2, a_3\,\&c.$
$a'_1, a'_2, a'_3\,\&c.\ a^2_1, a^2_2, a^2_3\,\&c. \ldots\ a^\beta_1, a^\beta_2, a^\beta_3\,\&c.$ sont des
coëfficiens arbitraires.

N.\textsuperscript{o} 5 si on \struck{remarque} remplace $\frac{gR^2}{2^2}$, $\frac{g'R^2}{2.3}$, $\frac{g''R^2}{2.5}$ \&c.
$\frac{g^\beta R^2}{2(n+1)}$ par $\theta$ on aura
\[ \uncertain{f}\theta = y = 1 - \frac{(n+1)\theta}{n+1} + \frac{(n+1)^2\theta^2}{1.2(n+1)(n+3)} - \frac{(n+1)^3\theta^3}{1.2.3(n+1)(n+3)(n+5)} \]
\[ + \frac{(n+1)^4\theta^4}{1.2.3.4(n+1)(n+3)(n+5)(n+7)} - \&c. \]
\[ \frac{dy}{d\theta} = -1 + \frac{(n+1)^2\theta}{(n+1)(n+3)} - \frac{(n+1)^3\theta^2}{1.2(n+1)(n+3)(n+5)} + \frac{(n+1)^4\theta^3}{1.2.3(n+1)(n+3)(n+5)(n+7)} - \&c. \]
\[ \frac{2\theta}{n+1}.\frac{d^2y}{d\theta^2} = \frac{(n+1)^{\uncertain{3}}.2\theta}{(n+1)(n+3)} - \frac{4(n+1)^2\theta^2}{1.2(n+1)(n+3)(n+5)} + \frac{6.(n+1)^3\theta^3}{1.2.3(n+1)(n+3)(n+5)(n+7)} - \&c. \]
et parconséquent
\[ y + \frac{dy}{d\theta} + \frac{2\theta}{n+1}.\frac{d^2y}{d\theta^2} = 0 \]
\[ \frac{dy}{d\theta} + \frac{n+3}{n+1}.\frac{d^2y}{d\theta^2} + \frac{2\theta}{n+1}.\frac{d^3y}{d\theta^3} = 0 \]
\[ \frac{d^2y}{d\theta^2} + \frac{n+5}{n+1}.\frac{d^3y}{d\theta^3} + \frac{2\theta}{n+1}.\frac{d^4y}{d\theta^4} = 0 \]
\&c.
\[ \frac{d^iy}{d\theta^i} + \frac{n+2i+1}{n+1}.\frac{d^{i+1}y}{d\theta^{i+1}} + \frac{2\theta}{n+1}\frac{d^{i+2}y}{d\theta^{i+2}} = 0 \]
En appliquant ici les remarques qui font l'objet de l'art
308. on verra qu'en donnant a $\theta$ une valeur positive
\note{En haut à droite, à l'encre, « 336 » ; au-dessus, un « 7 » barré
d'un trait oblique, sa pagination. Le premier « qui » est biffé, ce qui
fait de la liste de la première ligne le sujet de « désignent ». Ces lignes répètent, sans rature, celles qui sont barrées au dos de la
vue 678 (vue 679, lot 34), avec $A'$, $B'$ … et $B^0$. Le
premier des coëfficiens $B$ porte de nouveau un exposant noyé dans un
pâté, comme à la vue 682 ; « 0 » reste une lecture offerte. « N.o 5 »
est écrit en tête de l'alinéa : c'est son numéro d'article, auquel
renvoie « (art 5) » à la vue 686, comme « art 4 » (vue 688) renvoie à
l'intégrale de la vue 682. Les dénominateurs « $2^2$ »,
« $2.3$ », « $2.5$ » sont lus tels qu'écrits, là où la règle qu'elle
énonce ensuite, $2(n+1)$ pour $n = 1, 3, 5$, donnerait $2.2$, $2.4$, $2.6$ ;
on garde la feuille. Dans la marge de gauche, devant « $y =$ », une
lettre à longue hampe suivie de $\theta$ et d'un signe égal, lue avec
doute $f\theta$ ; le $y$ est récrit sur un premier tracé. Dans la
troisième équation, l'exposant « 3 » de $(n+1)$ est très pâle, comme
effacé ; le calcul demande $(n+1)$ sans exposant. Son 4 est écrit comme
un « b » (« $1.2.3.4$ », « $4(n+1)^2\theta^2$ »). « l'art 308. » : le
mot « art » est lu d'après le tracé, chargé d'encre. La phrase se
poursuit à la vue 686 ; le bas du feuillet est blanc.}

\page{686}

qui rende nulle la fluxion $\frac{d^{i+1}y}{d\theta^{i+1}}$ \struck{une valeur positive qui}
les deux autres termes de cette équation recevront des valeurs
opposées et on en conclura que les differentes valeurs de
$g$, $g'$, $g^2 \ldots g^\beta$ sont toutes réelles.

N.\textsuperscript{o} 6 Examinons à présent la forme du facteur
$A^\beta\,\text{sin}.\beta\varphi + B^\beta\,\text{cos}.\beta\varphi$ qui est dû a la difference de temperature
entre les divers points d'une même couche: ce facteur
contient les deux quantités angulaires $\text{sin}.\beta\varphi$ et $\text{cos}.\beta\varphi$ et
parconséquent, soit qu'on suppose $\beta$ infini, ou nul, il
doit se réduire a un seul terme. La première supposition
\uncertain{indiqueroit} que sur la circonférence d'\struck{es différentes couches} \add{de chacune des couches}
il y \struck{auroit} \add{existeroit} un nombre infini de points qui auroient
\struck{une} \struck{seule et} \add{la} même temperature, \struck{La secon} elle donneroit lieu
\marginal{a la condition} $\text{cos}\,\beta\varphi = 0$. La seconde exprime qu'a raison de l'uniformité
dela chaleur dans \struck{les chacune} \add{tous les points. qui sont} \struck{également distants} \add{a une distance égale} \add{de l'axe du cylindre} des couches, \struck{\uncertain{la}} $v$ n'est
pas fonction delaquantité angulaire $\varphi$; \struck{elle donne} \add{on en conclut}
\marginal{l'équation} $\text{sin}\,\beta\varphi = 0$.

Les deux suppositions dont on vient de parler \struck{ne} sont
donc \struck{que} deux manières differentes de présenter \struck{un seul}
\struck{et} le même état dela distribution dela chaleur.

Il resulte delà que plus le nombre $\beta$ est grand
plus aussi les differentes valeurs de $g^\beta$ doivent être grandes,
car s'il en étoit autrement la supposition de $\beta$ infini
réduiroit la valeur de $s$, \add{(art 5)} a l'unité, et on seroit mené
à conclure que la supposition de l'uniformité dela
chaleur entre les differens points d'une même couche
\note{En haut à droite, à l'encre, « 337 » ; à sa gauche, un « 8 » barré
d'un trait oblique, sa pagination. La première ligne continue la phrase
de la vue 684 (« en donnant a $\theta$ une valeur positive | qui rende
nulle la fluxion… ») ; les mots « une valeur positive qui », après la
fraction, sont barrés d'un trait qui les coupe à mi-hauteur. « indiqueroit »
est écrit « indequeroit », l'i sans point et formé comme un e. Dans
« d'es différentes couches », l'« es » et « différentes » sont barrés ;
« de chacune des couches » est écrit au-dessus, plus fin ; « couches »,
sur la ligne, n'est pas barré. « existeroit » est écrit au-dessus
d'« auroit » barré, « la » au-dessus de « seule et » barré. « a la
condition » et « l'équation » sont écrits dans la marge de gauche, plus
petit, devant les deux égalités qu'ils introduisent. L'addition
au-dessus de « les chacune des couches » se lit « tous les points. qui
sont également distants de l'axe du cylindre », « également distants »
barré et « a une distance égale » écrit plus haut encore, au-dessus ;
aucun signe n'en marque la place, et « des couches » n'est pas barré sur
la ligne. Avant $v$, un petit mot barré, lu avec doute « la ». « on en
conclut » est écrit en fin de ligne, dans l'interligne, au-dessus de
« elle donne » barré. « (art 5) » est écrit au-dessus de « $s$, ». La
phrase se poursuit à la vue 688 ; le bas du feuillet est blanc.}

\page{688}

entraine nécessairement la même uniformité entre tous les
points du cylindre.

\struck{En admettant donc} \add{si on admet} qu'en suivant l'ordre des indices
$1, 2 \ldots \beta$ les valeurs des quantités $g\ g' \ldots g^\beta$ sont \struck{soient}
toutes plus grandes que celles qui appartiennent a l'ordre
précédent et qu'en même tems, elles soient plus petites que
les valeurs qui conviennent a l'ordre suivant.
On verra \add{donc} que plus le tems écoulé augmente et plus les \add{\ill{}}
termes qui composent la valeur de $\underline{v}$ dans l'équation
trouvée art 4 deviennent petits par rapport à ceux
qui les précèdent.

Il y a donc une certaine valeur de $\underline{t}$ pour laquelle
le mouvement dela chaleur commence a être \uncertain{sensiblement}
représenté par l'équation
\[ v = B\left\{a_1e^{-g_1kt}s + a_2e^{-g_2kt}s + \&c\right\} \]
\[ + (A'\,\text{sin}.\varphi + B'\,\text{cos}\,\varphi)\left\{a'_1e^{-g'_1kt}s + a'_2e^{-g'_2kt}s + \&c\right\} \]
Si dans cet état on choisit deux points appartenans
a une même couche et situés aux extrémités d'un même
diamètre; en représentant par $\varphi_1$ et $\varphi_2$ leurs distances
respectives à l'origine, par $v_1$ $v_2$, leurs temperatures
correspondantes au tems $\underline{t}$, on \struck{aura} \struck{verra} \add{observera} que la difference
entre la valeur de $v_1$ et celle de $v_2$ depend uniquement
de \struck{la difference} \add{celle qui existe} entre \struck{celle} les facteurs $(A'\,\text{sin}\,\varphi_1 + B'\,\text{cos}\,\varphi_1)$
et $(A'\,\text{sin}.\varphi_2 + B'\,\text{cos}.\varphi_2)$ et qu'a cause de $\varphi_2 = \pi + \varphi_1$
la somme de ces deux facteurs est nulle.
\note{En haut à droite, à l'encre, « 338 » ; à sa gauche, un « 9 » barré
d'un trait oblique, sa pagination. La première ligne achève la phrase de
la vue 686. « si on admet » est écrit au-dessus de « En admettant donc »
barré ; « soient » après « sont » est barré, puis récrit sur la ligne
suivante, chargé d'encre. « donc » est écrit au-dessus de « verra » ;
au-dessus de « les », en fin de ligne, un petit mot ajouté, au bord du
feuillet, qui n'a pas pu être lu. « sensiblement » est offert pour un
mot rapide dont le tracé se lit à peu près « stablement ». Dans la
formule, $B$ n'a pas d'exposant et les $s$ n'ont pas d'indice. « la
valeur » : l'article est repassé. Au-dessus de « la difference » barré,
« celle qui existe » ; au-dessus de « les », un « celle » écrit puis
barré. « observera » est écrit au-dessus d'« aura » et de « verra »,
barrés l'un et l'autre. Le bas du feuillet est blanc ; la vue 689 en est
le dos, blanc.}

\page{690}

parconséquent la somme des seconds termes des valeurs de
$v_1$ et $v_2$ est également nulle et on a
\[ \frac{v_1+v_2}{2} = B\left(a_1e^{-g_1kt}s + a_2e^{-g_2kt}s + \&c\right) \]
Ainsi la demi-somme des temperature des points opposés est égale
a la temperature moyenne dela couche a laquelle ils appartiennent,
si on veut que $v_1$ et $v_2$ représentent le maximum et le minimum
des temperatures existantes dans la même couche on aura
$\frac{dv_1}{d\varphi} = 0$ $\frac{dv_2}{d\varphi} = 0$ et parconséquent $A'\,\text{cos}\,\varphi_1 - B'\,\text{sin}\,\varphi_1 = 0$, $\frac{A'}{B'} = \text{tang}.\varphi_1$,
$v'_1$ et $v'_2$ étant les temperatures de deux autres points placés a la
distance $\frac{\pi}{2}$ des premiers, les facteurs angulaires qui multiplient
les seconds termes des valeurs de $v'_1$ et $v'_2$ seront
\[ B'(\text{tang}.\varphi_1\,\text{cos}\,\varphi_1 - \text{sin}\,\varphi_1 = 0 \quad\text{et}\quad - B'(\text{tang}\,\varphi_1\,\text{cos}\,\varphi_1 - \text{sin}\,\varphi_1) = 0 \]
Les valeurs de $v'_1$ et $v'_2$ se réduiront \add{donc} au seul terme qui
exprime la temperature moyenne dela couche.

Ainsi lorsqu'il n'y a aucun mouvement de la chaleur
dans le sens de l'axe du cylindre, quelque soit d'ailleurs
la manière dont ce solide ait été échauffé il doit arriver
lorsque le tems $\underline{t}$ aura acquis une certaine valeur,
que le refroidissement s'opere, dans les points situés
sur le rayon qui partage en deux parties égales l'arc
compris entre les points qui ont la plus grande et la
moindre chaleur; absolument dela même manière que si
ce rayon appartenoit au cylindre dont tous les point
egalement eloignés del'axe \struck{ont} \add{auroient} reçus une temperature
uniforme.

Il existe la plus grande analogie entre cequi arrive ici et ce
qui a lieu dans l'armille; aussi pourroit-on appliquer au
cas présent tout cequ'on lit art 245 p 277 dela Théorie dela chaleur
On sentira la raison de cette identité des lois dela chaleur
\note{En haut à droite, à l'encre, « 339 » ; au-dessus, « 10 » barré d'un
trait oblique, sa pagination. Le premier mot continue la vue 688.
Dans la première des deux égalités du bas, la parenthèse ouverte n'est
pas fermée ; on garde la feuille. « donc » est écrit au-dessus de
« réduiront », « auroient » au-dessus d'« ont » barré. « 245 » : le 4
en forme de « b », comme ailleurs dans sa main ; « p 277 » lu tel
qu'écrit. « Théorie dela chaleur » n'est pas souligné. La phrase se
poursuit au dos du feuillet, vue 691, dont l'écriture se voit ici à
l'envers à travers le papier ; en haut à gauche, deux petits traits
obliques pâles, qui ne sont pas un nombre.}

\page{691}

si on remarque que chacune des tranches parallèles a la base
circulaire du cylindre doit se comporter comme une plaque
ronde qui seroit soustraite au refroidissement dû a l'action
del'air et que le cas del'armille est correspond
a celui ou on feroit le rayon constant.
\note{Dos du feuillet de la vue 690, écrit sur ses cinq premières lignes
seulement ; le reste est blanc. Ni pagination ni nombre à l'encre : ce
qu'on voit pâle en haut à gauche est le « 10 » et le « 339 » du recto
vus à travers le papier. « est correspond » : ainsi sur la feuille, sans
biffure visible. La phrase achève celle de la vue 690 (« On sentira la
raison de cette identité des lois dela chaleur | si on remarque… ») ; le
feuillet suivant (vue 692) ouvre un nouveau titre.}

\page{692}

Equation du mouvement varié dela chaleur dans un cylindre
solide

1. Supposons que toutes les tranches parallèles à la base d'un
cylindre solide d'une longueur infinie et dont le côté est
perpendiculaire à la base circulaire, ayant été échauffé d'une
manière entièrement arbitraire, mais semblable pour chacune
d'elles; le cylindre soit exposé ensuite à un courant d'air
plus froid.

Il s'agit de trouver les équations qui doivent représenter le
mouvement dela chaleur dans un tel solide.

On voit d'abord que l'équation générale
\[ \frac{d\uncertain{v}}{dt} = \frac{K}{C.D}\left(\frac{d^2v}{dx^2} + \frac{d^2v}{dy^2} + \frac{d^2v}{dz^2}\right) \]
qui sert à déterminer le mouvement dela chaleur
dans l'intérieur des corps solides se réduit ici à
\[ \frac{dv}{dt} = \frac{K}{C.D}\left(\frac{d^2v}{dx^2} + \frac{d^2v}{dy^2}\right) \ldots\ldots (A) \]
car, suivant l'hypothèse, les tranches du cylindre étant
échauffés d'une manière semblable, tous les points qui
sont situés sur des lignes parallèles a son axe, qui est
en même tems celui des $z$, ont la même temperature,
parconséquent le mouvement dela chaleur est nulle dans
cette direction et on a $\frac{d^2v}{dz^2} = 0$.

Par la même raison l'équation générale relative a la surface
se réduit à
\[ m\frac{dv}{dx} + n\frac{dv}{dy} + \frac{h}{K}vq = 0 \ldots\ldots (B) \]

2. Le centre d'une des tranches parallèles à la base du cylindre
étant pris pour origine; soit $\underline{r}$ le rayon du cercle qui mesure
une des couches de cette tranche et $\varphi$ l'arc compris
\note{Le titre et les articles 1 et 2 du morceau qui commence à la vue
672 (lot 34), écrits une seconde fois presque mot pour mot, d'une main
posée, en tête d'un feuillet sans pagination de sa main ; en haut à droite, à l'encre, « 340 ». Sous le
titre, un court trait. Les numéros d'article « 1. » et « 2. » sont dans
la marge. Au numérateur du premier membre de la première équation, la
lettre après d est chargée d'encre et peut se lire x ; $v$ est offert
d'après l'équation (A). L'accord de « échauffé » et d'« échauffés » est
celui de la feuille. La phrase se poursuit au dos, vue 693 ; à travers
le papier se voit à l'envers l'écriture de ce dos.}

\page{693}

entre un des points de la circonference et celui auquel appartient
l'ordonnée $y = 0$, on aura quelque soit $\varphi$. $x = r\,\text{cos}\,\varphi$, $y = r\,\text{sin}.\varphi$;
d'ou on tire: $dx = dr\,\text{cos}\,\varphi - r\,d\varphi\,\text{sin}\,\varphi$, $dy = dr\,\text{sin}.\varphi + r\,d\varphi\,\text{cos}.\varphi$,
\[ dr = dx\,\text{cos}.\varphi + dy\,\text{sin}.\varphi, \qquad d\varphi = \frac{dy\,\text{cos}\,\varphi - dx\,\text{sin}.\varphi}{r} \]
\[ \frac{dr}{dx} = \text{cos}.\varphi, \quad \frac{dr}{dy} = \text{sin}\,\varphi; \quad \frac{d\varphi}{dx} = -\frac{\text{sin}\,\varphi}{r}, \quad \frac{d\varphi}{dy} = \frac{\text{cos}\,\varphi}{r} \]
ces valeurs de $\frac{dr}{dx}$, $\frac{dr}{dy}$, $\frac{d\varphi}{dx}$, et $\frac{d\varphi}{dy}$ étant mises dans celles de
$\frac{dv}{dx}$, $\frac{dv}{dy}$, $\frac{d^2v}{dx^2}$ et $\frac{d^2v}{dy^2}$ on trouve:
\[ \frac{dv}{dx} = \text{cos}.\varphi\,\frac{dv}{dr} - \frac{\text{sin}.\varphi}{r}\,\frac{dv}{d\varphi} \]
\[ \frac{dv}{dy} = \text{sin}.\varphi\,\frac{dv}{dr} + \frac{\text{cos}.\varphi}{r}\,\frac{dv}{d\varphi} \]
\[ \frac{d^2v}{dx^2} = \frac{\text{sin}^2\varphi}{r}\,\frac{dv}{dr} + \frac{2\,\text{sin}.\varphi\,\text{cos}.\varphi}{r^2}.\frac{dv}{d\varphi} + \text{cos}^2\varphi.\frac{d^2v}{dr^2} - \frac{2\,\text{sin}\,\varphi\,\text{cos}\,\varphi}{r}.\frac{d^2v}{dr\,d\varphi} + \frac{\text{sin}^2\varphi}{r^2}.\frac{d^2v}{d\varphi^2} \]
\[ \frac{d^2v}{dy^2} = \frac{\text{cos}^2\varphi}{r}.\frac{dv}{dr} - \frac{2\,\text{sin}.\varphi\,\text{cos}.\varphi}{r^2}.\frac{dv}{d\varphi} + \text{sin}^2\varphi\,\frac{d^2v}{dr^2} + \frac{2\,\text{sin}.\varphi\,\text{cos}.\varphi}{r}\,\frac{d^2v}{dr\,d\varphi} + \frac{\text{cos}^2\varphi}{r^2}\,\frac{d^2v}{d\varphi^2} \]
\[ \frac{d^2v}{dx^2} + \frac{d^2v}{dy^2} = \frac{1}{r}\,\frac{dv}{dr} + \frac{d^2v}{dr^2} + \frac{1}{r^2}\,\frac{d^2v}{d\varphi^2} \]
et l'équation (A) devient
\[ \frac{dv}{dt} = \frac{K}{C.D}\left\{\frac{1}{r}\,\frac{dv}{dr} + \frac{d^2v}{dr^2} + \frac{1}{r^2}\,\frac{d^2v}{d\varphi^2}\right\} \ldots\ldots (C) \]
A l'egard de l'équation relative a la surface, il faut remarquer
que la condition $m\,dx + n\,dy = 0$, qui donne les valeurs de $\underline{m}$ et de $\underline{n}$,
et parconséquent aussi celle de $\underline{q}$ étant, ici, $x\,dx + y\,dy = 0$; on
peut mettre aulieu de $\underline{m}$, $\underline{n}$ et $q$; $x$, $y$ et $\sqrt{x^2+y^2}$, ou plutôt leurs valeurs
$r\,\text{cos}\,\varphi$, $r\,\text{sin}.\varphi$ et $r$. Si on met aussi a la place de $\frac{dv}{dx}$ et $\frac{dv}{dy}$
les valeurs trouvées plus haut l'équation (B) deviendra:
\[ r\,\text{cos}.\varphi\left(\text{cos}\,\varphi\,\frac{dv}{dr} - \frac{\text{sin}.\varphi}{r}\,\frac{dv}{d\varphi}\right) + r\,\text{sin}.\varphi\left(\text{sin}.\varphi\,\frac{dv}{dr} + \frac{\text{cos}.\varphi}{r}.\frac{dv}{d\varphi}\right) + r\frac{h}{K}v = 0 \]
elle se réduira donc à
\[ \frac{dv}{dr} + \frac{h}{K}v = 0 \ldots\ldots (D) \]

3. En traitant la question d'une manière spéciale on obtiendroit
egalement les deux équations (C) et (D). Pour y parvenir il
\note{Dos du feuillet de la vue 692, sans nombre à l'encre ; en haut à
droite, « 2 » souligné, sa pagination de ce morceau (la page qui le
précède, vue 692, n'en porte pas). La première ligne continue la
phrase de la vue 692. Dans la liste « $\frac{d^2v}{dy^2}$ » de la
cinquième ligne, la lettre du numérateur est récrite et se lit aussi
$x$. Dans « les valeurs trouvées plus haut », l'article est récrit sur
un premier mot. Le numéro d'article « 3. » est dans la marge. Le texte est celui des
vues 674 et 676 (lot 34), sans la ligne barrée ni les ratures. La page
s'arrête sur « Pour y parvenir il » ; la vue 694 ne la continue pas
(voir sa note).}

\page{694}

manières differentes de présenter le même état dela distribution dela
chaleur. Il résulte delà que plus le nombre $\beta$ est grand plus aussi
les differentes valeurs de $g^\beta$ doivent être grandes, car, s'il en
étoit autrement la supposition de $\beta$ infini réduiroit la valeur
de $s$ (art 5) a l'unité, et on seroit forcé de conclure que la
supposition de l'uniformité dela chaleur entre les differens points
d'une même couche, entraîne nécessairement la même uniformité
entre tous les points du cylindre.

Si on admet qu'en suivant l'ordre des indices $1, 2, \ldots \beta$ les valeurs
des quantités $g$, $g'$, $g^2 \ldots g^\beta$, soient toutes plus grandes que celles
qui appartiennent a l'ordre précédent et qu'en mêmetems, elles
soient plus petites que les valeurs qui conviennent a l'ordre
suivant, on verra donc que plus le tems écoulé augmente
et plus aussi les termes qui composent la valeur de $\underline{v}$ dans
l'équation trouvée art 4 deviennent petits par rapport a ceux
qui les précèdent. Ainsi il y a une certaine valeur de $\underline{t}$
pour laquelle le mouvement dela chaleur commence à être
sensiblement représenté par l'équation:
\[ v = B\left\{a_1e^{-g_1kt}s_1 + a_2e^{-g_2kt}s_2 + \&c\right\} \]
\[ + (A'\,\text{sin}.\varphi + B'\,\text{cos}.\varphi)\left\{a'_1e^{-g'_1kt}s'_1 + a'_2e^{-g'_2kt}s'_2 + \&c\right\}. \]
Dans cet état, si on choisit deux points appartenant à une
même couche et situés aux extrémités d'un même diamètre,
en représentant par $\varphi_1$ et $\varphi_2$ leurs distances respectives à l'origine,
par $v_1$ $v_2$ leurs temperature correspondantes au tems $\underline{t}$; on
observera que la difference entre la valeur de $v_1$ et celle de $v_2$
depend uniquement de celle qui existe entre les facteurs
$A'\,\text{sin}\,\varphi_1 + B'\,\text{cos}\,\varphi_1$ et $A'\,\text{sin}\,\varphi_2 + B'\,\text{cos}.\varphi_2$. A cause de $\varphi_2 = \pi + \varphi_1$
la somme de ces deux facteurs est nulle; la somme des
\note{En haut à droite, à l'encre, « 341 » ; à sa droite, un « 7 » barré
d'un trait oblique, sa pagination. Ce feuillet ne suit pas la vue 693,
qui s'arrête sur « Pour y parvenir il » : il reprend, mis au net en
une écriture plus régulière et presque sans rature, le texte des vues
686 et 688 (« manières differentes de présenter le même état… » jusqu'à
« la somme de ces deux facteurs est nulle »), avec des variantes de
détail (« forcé » pour « mené », « Ainsi il y a » pour « Il y a donc »,
les indices des $s$ écrits). Les pages 3 à 6 de cette mise au net ne
sont pas dans ce lot. Dans « au tems », « au » est traversé d'un trait
qui peut être son trait de liaison. L'exposant du second $A$ de la
dernière formule est plus haut que l'accent du premier, et lu $A'$.
Le bas du feuillet est blanc ; la phrase se poursuit au dos, vue 695.}

\page{695}

seconds termes des valeurs de $v_1$ et $v_2$ est donc également nulle et on a:
\[ \frac{v_1+v_2}{2} = B\left(a_1e^{-g_1kt}s_1 + a_2e^{-g_2kt}s_2 + \&c\right) \]
Ainsi la demi somme des temperatures des points opposés est
égale a la temperature moyenne dela couche a laquelle ils
appartiennent.

Si on veut que $v_1$ et $v_2$ représentent le maximum et le minimum
existantes dans la même couche on aura $\frac{dv_1}{d\varphi} = 0$, $\frac{dv_2}{d\varphi} = 0$ ou, cequi
est la même chose, $A'\,\text{cos}\,\varphi - A'\,\text{sin}\,\varphi = 0$, $\frac{A'}{B'} = \text{tang}.\varphi$. $v'_1$ et $v'_2$
étant les temperatures de deux autres points placés a la distance
$\frac{\pi}{2}$ des premiers, les facteurs angulaires qui multiplient les seconds
termes des valeurs de $v'_1$ et $v'_2$ se réduiront à $B'(\text{tang}.\varphi_1\,\text{cos}\,\varphi_1 - \text{sin}.\varphi_1)$
et $-B'(\text{tang}.\varphi_1\,\text{cos}\,\varphi_1 - \text{sin}.\varphi_1)$. Les valeurs de $v'_1$ et $v'_2$ seront
donc alors composées du seul terme qui exprime la temperature
moyenne dela couche.

Ainsi lorsqu'il n'y a aucun mouvement dela chaleur dans le sens
de l'axe du cylindre, quelque soit d'ailleurs la manière dont ce
solide ait été échauffé, il doit arriver, lorsque le tems $\underline{t}$
aura acquis une certaine valeur, que le refroidissement s'opère,
dans les points situés sur le rayon qui partage en deux parties
égales l'arc compris entre les points qui ont la plus grande
et la moindre chaleur, absolument de la même manière
que si ce rayon appartenoit au cylindre dont tous les points également
éloignés de l'axe auroient reçus une temperature uniforme.

Il existe la plus grande analogie entre cequi arrive ici et cequi a lieu
dans l'armille; aussi pourroit-on appliquer au cas présent tout
cequ'on lit art 245. On sentira la raison de cette identité des
lois de la chaleur si on remarque que chacune des tranches parallèles à
la base circulaire du cylindre doit se comporter comme une plaque
ronde qui seroit soustraite au refroidissement dû a l'action de l'air, et que
le cas de l'armille correspond à celui ou on feroit le rayon constant.
\note{Dos du feuillet de la vue 694, sans nombre à l'encre ; en haut à
droite, « 8 », non barré, sa pagination. C'est la mise au net des vues
690 et 691, qu'elle suit presque mot pour mot jusqu'à la même dernière
phrase ; le renvoi « art 245 » y est donné sans la page ni le titre de
l'ouvrage, et « est correspond » y devient « correspond ». « $A'\,\text{cos}\,\varphi
- A'\,\text{sin}\,\varphi = 0$ » : la seconde lettre est écrite $A'$ sur la
feuille, où le brouillon (vue 690) porte $B'$ ; on garde la feuille. Sous
« à celui », un gros trait d'encre horizontal, sans rapport apparent avec
le texte. Le texte s'arrête sur « constant. » ; le bas du feuillet est
blanc.}

\page{696}

Equations du mouvement dela chaleur \add{dans} la sphère, \struck{solide} le cône
et le cylindre\struck{s} solides échauffés d'une manière entièrement
arbitraire\struck{s}

Lorsqu'on veut avoir l'équation \add{relative} \struck{du mouvement dela chaleur dans} a la
\struck{une} sphère \struck{arbitraire} \add{\struck{échauffée d'une manière}}, on peut se contenter comme l'ont fait
M\textsuperscript{rs} dela Place et Poisson, de \struck{transformer} \add{\struck{substituer dans}} \add{transformer} l'équation \struck{du mouvement}
générale du mouvement dela chaleur dans un corps solide, en y
\struck{substituant} \add{\struck{substituant}} \add{mettant} \add{a la place} \struck{dans cette équation des valeurs des coordonnées}
\struck{propres a la figure} des coordonnées rectangulaires d'autres coordonnées
appropriées a la figure du corps que l'on considère.

Des transformations \struck{du même} analogues donne\uncertain{nt} également les
équations du mouvement dela chaleur dans \struck{le} cône et \add{dans} le cylindre.
mais \struck{il sera difficile}, en suivant ce procédé, \add{il seroit} d'appercevoir
la liaison qui existe entre les trois équations. \ill{}

\marginal{Au contraire} si \add{lorsque} \struck{on en suivant} conformement a la méthode employée
\struck{Dans la \uncertain{tous le}} par M\textsuperscript{r} fourier; on \struck{considère} \add{\struck{envisage}} \add{examine} chacun \struck{des} de
\struck{cas} ces questions d'une manière spéciale \struck{et qu'}on \struck{détermine}
\struck{verra que la molécule dela sphère dont on a} \add{\struck{dont}}, pour \struck{trouver}
\struck{l'équation relative a ce solide il faut connoître}
\struck{et par l'examen} \struck{et on} \add{\uncertain{qu'on cherche la \ill{} qu'elle et}} détermin\uncertain{e} \struck{\ill{}} la quantité de chaleur
acquise \struck{durant} \struck{perd} dans \struck{l'élément du} \add{un tems donné} tems, par une \struck{des}
molécules \struck{qui} \struck{appartiennent a chacun de ces solides} \add{du solide} on s'apperçoit
aisément que la molécule de la sphère appartient en
même tems a un solide conique qui \add{a} \struck{auroit} \struck{son} sommet
au centre dela sphère et qui \struck{seroit} \add{est} terminé par une
calotte sphérique. \struck{L'élément de ce} \struck{La molécule dont il
s'agit sera} \add{est} \struck{comprise d'une part} entre deux \struck{couches concentriques}
\struck{qui} calottes concentriques \struck{elle} \struck{l'est} \add{est terminée d'ailleurs} \struck{ses autres faces sont}
\struck{terminées par des plans} \struck{lorsqu'on veut avoir l'élément}

La seule difference qui puisse exister entre la molécule dont
il s'agit ici et celle qui appartiendroit au cône \struck{solide}
qui ayant le même centre \struck{et le} même axe et le même coté
\note{Feuillet de brouillon, très raturé ; en haut à droite, à l'encre,
« 342 », et au-dessus un « 1 » barré d'un trait oblique, sa pagination
d'un nouveau morceau. Le titre est suivi d'un court trait. L's final de
« cylindres » et d'« arbitraires » est annulé d'une petite croix : on
donne le singulier. Dans les lignes 4 à 10, les additions sont écrites
dans l'interligne au-dessus des mots barrés ; « transformer » est
récrit sous la ligne et relié à sa place par un trait courbe ; « a la
place » est écrit à côté de « mettant », au-dessus de « substituant »
barré. Un long trait de liaison court sous « Lorsqu'on veut avoir
l'équation » et sous « en y » sans être une biffure : ces mots sont
gardés. « donne\uncertain{nt} » : la fin du mot est récrite en surcharge
sur une première terminaison (« donneroient » ?), noyée d'encre.
« il seroit » est écrit au-dessus de « procédé, d'appercevoir » ;
« difficile », barré avec « il sera », n'est pas rétabli. En fin de la
ligne « les trois équations. », un mot commencé, non lu. « Au
contraire » est dans la marge de gauche, devant « si » ; « lorsque »
au-dessus d'« on en suivant » barré. Au-dessus de « considère » barré,
« envisage » barré, et au-dessous « examine ». Le passage de « verra
que la molécule… » à « par l'examen » est barré ligne à ligne ; les
mots écrits dans l'interligne au-dessus de « détermine », d'une plume
fine et serrée, ne sont lus qu'en partie, et un mot barré sous une tache
d'encre suit « détermin- ». Du passage « L'élément de ce… » à « avoir
l'élément », il n'est pas sûr, pour quelques mots (« entre deux »,
« calottes concentriques »), si le trait qui les traverse est une
biffure ou son trait de liaison ; on donne ce que la main semble garder.
Le bas du feuillet est blanc sous « coté » ; le dos (vue 697) est blanc.}

\page{698}

Lorsque la sphère est échauffée uniformement du centre a la
circonference les deux termes $\frac{d^2v}{d\theta^2}.\frac{1}{r^2}$; $\frac{d^2v}{d\varphi^2}\frac{1}{r^2\,\text{sin}^2\theta}$
sont l'un et l'autre nuls. Il est évident qu'il n'existe
alors aucun mouvement dela chaleur dans une invelappe
sphérique correspondante a une des valeur du rayon $\underline{r}$ comprise
entre zéro et la valeur entière du rayon dela sphère
car alors \add{tous} les differens quadrilatères que l'on peut former
sur la surface sphérique en menant du pôle les differens
méridiens qui coupent alafois l'équateur et \struck{tous} les parallèles
ont tous lamême temperature ainsi il ne passe aucune
quantité de chaleur par les côtés qui forment chacun
de ces quadrilatères.

L'équation $\frac{d^2v}{d\theta^2}.\frac{1}{r^2} + \frac{d^2v}{d\varphi^2}.\frac{1}{r^2\,\text{sin}^2\theta} = 0$ exprime la condition
dela constance des temperatures des differens quadrilatères
qui appartiennent \struck{aux} la surface sphérique correspondante
a la valeur $\underline{r}$ du rayon. Elle a lieu lorsqu'en prenant
\struck{le} la parallèle dont le rayon est $r\,\text{sin}\,\theta$, $\theta$ étant l'arc du
meridien compris entre le pôle et ce parallèle la quantité
de chaleur qui passe par l'arc \struck{$d\theta$ du meridien} $r\,\text{sin}\,\theta\,d\theta$
du \struck{paral} méridien, arc qui forme un de côté du quadrilatère,
est précisement egale a la quantité de chaleur qui pendant
le même tems $\frac{d\uncertain{v}}{dt}$, passe par l'arc $r\,d\varphi$ du paralelle,
arc qui forme le coté contigu du premier dans le quadrilatère
et qui appartient au parallèle dont la distance au pôle
est mesurée par l'arc $\theta - d\theta$.
\note{En haut à droite, à l'encre, « 343 » ; pas de pagination de sa
main. « invelappe » : ainsi. Dans « lamême », un pâté sur la fin du
mot. « \struck{aux} » est barré d'un trait appuyé ; « la surface » suit
sans correction. Après « l'arc », « $d\theta$ du meridien » est barré et
remplacé à la suite, sur la ligne, par « $r\,\text{sin}\,\theta\,d\theta$ » ;
le calcul attendrait $r\,d\theta$ pour l'arc du méridien et
$r\,\text{sin}\,\theta\,d\varphi$ pour celui du parallèle : on garde la
feuille. Au-dessus de « méridien », « paral » est écrit puis barré.
Dans « le même tems », la fraction qui suit a au numérateur, après d,
une petite lettre bouclée lue avec doute $v$ (elle peut être $\theta$). « paralelle »
et « parallèle » alternent sur la feuille. Le texte s'arrête sur
« $\theta - d\theta$: » ; le bas du feuillet est blanc, le dos (vue 699)
aussi.}

\end{document}
